% Atur variabel berikut sesuai namanya

% nama
\newcommand{\name}{Bagas Surya Wirawan}
\newcommand{\authorname}{Wirawan, Bagas Surya}
\newcommand{\nickname}{Bagas}
\newcommand{\advisor}{Dr. Ir. Djoko Purwanto, M.Eng.}
\newcommand{\coadvisor}{Fajar Budiman, S.T., M.Sc.}
\newcommand{\examinerone}{Ir. Tasripan, M.T.}
\newcommand{\examinertwo}{Ir. Harris Pirngadi, M.T.}
\newcommand{\examinerthree}{Alan Turing, ST., MT}
\newcommand{\headofdepartment}{Ronny Mardiyanto, S.T., M.T., Ph.D.}

% identitas
\newcommand{\nrp}{5022221026}
\newcommand{\advisornip}{196512111990021002}
\newcommand{\coadvisornip}{198607072014041001}
\newcommand{\examineronenip}{196204181990031004}
\newcommand{\examinertwonip}{196205101989031001}
\newcommand{\examinerthreenip}{18560710 194301 1 001}
\newcommand{\headofdepartmentnip}{18810313 196901 1 001}

% judul
\newcommand{\tatitle}{NAVIGASI ROBOT BERBASIS LAYERED COSTMAP UNTUK PERGERAKAN ANTAR LANTAI PADA LINGKUNGAN INDOOR PASCA GEMPA BUMI}
\newcommand{\engtatitle}{\emph{ROBOT NAVIGATION USING LAYERED COSTMAP FOR INTER-FLOOR MOVEMENT IN POST-EARTHQUAKE INDOOR ENVIRONMENTS}}

% tempat
\newcommand{\place}{Surabaya}

% jurusan
\newcommand{\studyprogram}{Teknik Elektro}
\newcommand{\engstudyprogram}{Electrical Engineering}

% fakultas
\newcommand{\faculty}{Teknologi Elektro dan Informatika Cerdas}
\newcommand{\engfaculty}{Intelligent Electrical and Informatics Technology}

% singkatan fakultas
\newcommand{\facultyshort}{FTEIC}
\newcommand{\engfacultyshort}{ELECTICS}

% departemen
\newcommand{\department}{Teknik Elektro}
\newcommand{\engdepartment}{Electrical Engineering}

% kode mata kuliah
\newcommand{\coursecode}{EE234899}
